\documentclass[letterpaper,11pt]{article}

\usepackage[letterpaper,left=0.80in,right=0.80in,top=0.70in,bottom=0.70in]{geometry}
\usepackage[hidelinks]{hyperref}
\usepackage[english]{babel}
\input{glyphtounicode}

\pagestyle{empty}
\setlength{\parindent}{0pt}
\setlength{\parskip}{10pt}
\pdfgentounicode=1

\begin{document}

{\Large\textbf{Aryan Miriyala}}\\
\href{mailto:aryanmiriyala@gmail.com}{aryanmiriyala@gmail.com} $|$ +1 419-315-0444 $|$
\href{https://www.linkedin.com/in/aryan-miriyala-7788a21b6/}{LinkedIn} $|$
\href{https://github.com/aryanmiriyala}{GitHub}

August 4, 2026

Pavago Hiring Team

Dear Pavago Hiring Team,

I am applying for the AI Product Engineer (Full Stack) role because the posting is centered on the kind of building I have kept choosing: taking an AI idea past a demo and making the surrounding product, backend, data model, and reliability checks strong enough for real users. I am especially interested in the ownership described here, from frontend and backend architecture to prompt handling, failure recovery, auth, permissions, and iteration from user feedback.

At SmartSolve, I architected a Next.js and TypeScript onboarding tracker with PostgreSQL, Drizzle ORM, SSO, and authentication middleware for sensitive internal workflows. I also scoped an AI-enabled QMS dependency register where the hard part was not only writing code, but turning an unclear operational problem into a system design that people could review and build from. I used Codex and Claude Code daily for planning, debugging, implementation review, and self-review, while keeping proprietary work isolated in Docker/devcontainer environments.

My projects map closely to Pavago's AI product needs. Fix-It-Flow is a Next.js PWA with typed API routes, camera input, voice commands, Gemini Vision, LLM reasoning, ElevenLabs audio, and DynamoDB-backed session persistence for repair guidance. RocketGrader used TypeScript/Express APIs, MongoDB models, Auth0 dashboards, S3 uploads, document parsing, LangChain, and Mistral AI to support AI-assisted assignment feedback. FalconGraph added the RAG side, with Python document processing, FAISS retrieval, FastAPI, Next.js, and cited answers that users could trace back to source material.

I would bring Pavago a full-stack AI builder profile across React/Next.js, TypeScript, Node.js, Python, REST APIs, PostgreSQL, AI APIs, RAG, authentication, Docker, Vercel, and AWS-backed workflows. More importantly, I would bring a practical approach to AI reliability: assume outputs can be wrong, design around edge cases, keep state and permissions clear, and make the product useful even when the model needs guardrails.

Sincerely,\\
Aryan Miriyala

\end{document}
